\chapter{Espérance et variables indicatrices}\label{ch:iv1:expectation-and-indicators}

\begin{ivtrains}
Calculer une moyenne sans jamais déterminer la loi~: décomposer un comptage en variables indicatrices, sommer
des probabilités de queue, utiliser la symétrie, puis aller un pas plus loin jusqu'à la variance du comptage. Le chapitre
commence par des échauffements rapides et se termine par des problèmes qui demandent une idée réelle qu'on doit trouver soi-même.
\end{ivtrains}

\begin{question}[$\star$ \iqroles{QR, QT, Trader, Analyste, Risque} \iqfirm{toutes} \iqlevel{junior}]\label{iq:iv1:expectation-and-indicators:1}
Deux dés équilibrés à six faces sont lancés, et vous gardez le plus grand chiffre affiché. En moyenne, que gardez-vous~?
\end{question}

\begin{question}[$\star$ \iqroles{QR, QT, Trader, Analyste} \iqfirm{toutes} \iqlevel{junior}]\label{iq:iv1:expectation-and-indicators:2}
Vous lancez un dé équilibré jusqu'à ce qu'un six apparaisse, et vous additionnez tous les chiffres obtenus, y compris ce dernier six. Quelle
est l'espérance du total~?
\end{question}

\begin{question}[$\star$ \iqroles{QR, QT, Trader, DS} \iqfirm{teneur de marché, hedge fund, crypto} \iqlevel{junior}]\label{iq:iv1:expectation-and-indicators:3}
Vous battez un jeu standard de 52 cartes et les étalez en une rangée. En moyenne, combien de paires de cartes voisines
ont la même valeur, comme deux rois côte à côte~?
\end{question}

\begin{question}[$\star$ \iqroles{QR, QT, Trader, Risque} \iqfirm{toutes} \iqlevel{junior}]\label{iq:iv1:expectation-and-indicators:4}
Deux jeux de 52 cartes sont battus séparément. Vous retournez une carte de chaque jeu en même temps, et vous poursuivez
à travers les deux jeux. Une coïncidence a lieu lorsque les deux cartes sont identiques en valeur et en couleur. Calculez l'espérance et la variance du nombre de coïncidences.
\end{question}

\begin{question}[$\star$ \iqroles{QR, QT, Trader, Analyste, Dev} \iqfirm{toutes} \iqlevel{junior}]\label{iq:iv1:expectation-and-indicators:5}
Vous lancez un dé équilibré jusqu'à ce que chacune des six faces soit apparue au moins une fois. Quelle est l'espérance du nombre de
lancers~?
\end{question}

\begin{question}[$\star$ \iqroles{QR, QT, Trader, DS} \iqfirm{toutes} \iqlevel{junior}]\label{iq:iv1:expectation-and-indicators:6}
Vous lancez une pièce équilibrée 21 fois. Une série est un bloc maximal de résultats identiques, de sorte que les lancers HHTHHH contiennent
trois séries. Quelle est l'espérance du nombre de séries dans vos 21 lancers, et quelle est sa variance~?
\end{question}

\begin{question}[$\star\star$ \iqroles{QR, QT, Trader} \iqfirm{toutes} \iqlevel{junior}]\label{iq:iv1:expectation-and-indicators:7}
Vous retournez un jeu de 52 cartes battu, une carte à la fois. En moyenne, à quelle position le deuxième as
apparaît~?
\end{question}

\begin{question}[$\star\star$ \iqroles{QR, QT, Trader, Risque} \iqfirm{toutes} \iqlevel{junior}]\label{iq:iv1:expectation-and-indicators:8}
Trois dés équilibrés dont les faces sont numérotées de 1 à 100 sont lancés. Quelle est l'espérance exacte du plus petit
des trois nombres obtenus~?
\end{question}

\begin{question}[$\star\star$ \iqroles{QR, QT, QD, Dev} \iqfirm{toutes} \iqlevel{junior}]\label{iq:iv1:expectation-and-indicators:9}
Vous lancez un dé équilibré jusqu'à ce qu'un chiffre sorte pour la deuxième fois. Quelle est l'espérance du nombre de lancers, sous forme de
fraction exacte~? Ensuite~: des clés sont placées dans une table de hachage de 1\,024 compartiments, chaque clé allant dans un compartiment
tiré uniformément au hasard. Environ combien de clés entrent avant la première collision~?
\end{question}

\begin{question}[$\star\star$ \iqroles{QR, QT, Trader} \iqfirm{toutes} \iqlevel{junior}]\label{iq:iv1:expectation-and-indicators:10}
Vous tirez des nombres aléatoires indépendants, chacun uniformément réparti entre 0 et 1, jusqu'à ce que leur somme dépasse 1. En moyenne, combien
de nombres tirez-vous~?
\end{question}

\begin{question}[$\star\star$ \iqroles{QR, QT, Risque, DS} \iqfirm{toutes} \iqlevel{junior, confirmé}]\label{iq:iv1:expectation-and-indicators:11}
Vous tirez un nombre d'une loi continue et vous le notez. Vous continuez ensuite à tirer des nombres indépendants
issus de la même loi jusqu'à ce que l'un d'eux dépasse le nombre noté. Combien de tirages supplémentaires
vous faut-il en moyenne~? Quelle est la probabilité qu'un seul tirage supplémentaire suffise~?
\end{question}

\begin{question}[$\star\star$ \iqroles{QR, QT, Risque} \iqfirm{toutes} \iqlevel{junior}]\label{iq:iv1:expectation-and-indicators:12}
100 nombres distincts sont écrits dans un ordre aléatoire et lus de gauche à droite. Un nombre est un record s'il est
plus grand que tous ceux qui le précèdent, le premier nombre étant compté comme un record. Calculez l'espérance et la variance du
nombre de records.
\end{question}

\begin{question}[$\star\star$ \iqroles{QR, QT, DS, Risque} \iqfirm{toutes} \iqlevel{junior}]\label{iq:iv1:expectation-and-indicators:13}
Vous lancez un dé équilibré 60 fois. Soient $A$ le nombre de 1 et $B$ le nombre de 2. Quelles sont la covariance
et la corrélation de $A$ et $B$~?
\end{question}

\begin{question}[$\star\star\star$ \iqroles{QR, QT, Trader} \iqfirm{teneur de marché, hedge fund, crypto} \iqlevel{junior, confirmé}]\label{iq:iv1:expectation-and-indicators:14}
Vous lancez un dé équilibré jusqu'à ce qu'un six apparaisse. Quelqu'un vous dit que chaque lancer effectué en chemin était pair, c'est-à-dire 2, 4
ou 6. Sachant cela, quelle est l'espérance du nombre de lancers~?
\end{question}

\begin{question}[$\star\star\star$ \iqroles{QR, QT, Trader} \iqfirm{teneur de marché, hedge fund} \iqlevel{junior, confirmé}]\label{iq:iv1:expectation-and-indicators:15}
Une marmite contient 100 nouilles cuites, ce qui fait 200 bouts libres. Vous prenez au hasard deux bouts libres et vous les
nouez ensemble, et vous recommencez jusqu'à ce qu'il ne reste plus aucun bout libre. Certaines des nouilles nouées se sont refermées en boucles. Quelle est
l'espérance du nombre de boucles~?
\end{question}

\begin{question}[$\star\star\star$ \iqroles{QR, QT, Trader, QD} \iqfirm{teneur de marché, crypto, hedge fund} \iqlevel{junior, confirmé}]\label{iq:iv1:expectation-and-indicators:16}
100 fourmis se tiennent à des positions indépendantes et uniformément aléatoires sur un bâton d'un mètre de long. Chacune est orientée vers la gauche ou vers la droite avec
une probabilité égale, et toutes avancent à un mètre par minute. Lorsque deux fourmis se rencontrent, elles font demi-tour instantanément. Une fourmi
qui atteint l'une ou l'autre extrémité tombe. En moyenne, combien de temps faut-il attendre jusqu'à ce que la dernière fourmi tombe~? Quelle est l'espérance du nombre
de collisions~?
\end{question}

\begin{question}[$\star\star\star$ \iqroles{QR, QT, Trader} \iqfirm{banque, teneur de marché, hedge fund} \iqlevel{junior, confirmé}]\label{iq:iv1:expectation-and-indicators:17}
Vous retournez un jeu de 52 cartes battu, une carte à la fois, et vous arrêtez dès que vous avez vu au moins un as et au
moins un roi. En moyenne, combien de cartes avez-vous retournées~?
\end{question}

\begin{question}[$\star\star\star$ \iqroles{QR, QT, Trader, DS} \iqfirm{toutes} \iqlevel{junior, confirmé}]\label{iq:iv1:expectation-and-indicators:18}
Un tiroir contient 5 paires de chaussettes, toutes différentes, soit 10 chaussettes au total. Vous sortez des chaussettes au hasard, une à la
fois, jusqu'à ce que vous ayez une paire assortie. Quelle est l'espérance du nombre de chaussettes que vous sortez~?
\end{question}

\begin{question}[$\star\star\star$ \iqroles{QR, QT, Trader} \iqfirm{toutes} \iqlevel{junior, confirmé}]\label{iq:iv1:expectation-and-indicators:19}
Une urne contient 6 boules rouges et 4 boules bleues. Vous tirez des boules au hasard sans remise et vous arrêtez dès qu'une couleur
est épuisée. En moyenne, combien de boules restent dans l'urne quand vous vous arrêtez~?
\end{question}

\begin{question}[$\star\star\star$ \iqroles{QR, DS, Risque, QD, Dev} \iqfirm{toutes} \iqlevel{junior, confirmé}]\label{iq:iv1:expectation-and-indicators:20}
Vous lancez 100 balles indépendamment dans 100 cases, chaque balle choisissant sa case uniformément au hasard. Calculez l'espérance et la variance du nombre de cases vides.
\end{question}

\begin{question}[$\star\star\star$ \iqroles{QR, QD, Dev, LLE} \iqfirm{toutes} \iqlevel{junior, confirmé}]\label{iq:iv1:expectation-and-indicators:21}
Quicksort reçoit 10 nombres distincts. Il choisit un pivot uniformément au hasard dans la liste courante, compare chaque
autre nombre de cette liste au pivot, sépare la liste entre les nombres plus petits et les plus grands, et trie chaque
partie de la même façon. Quelle est l'espérance du nombre de comparaisons~? Quelle est la formule pour $n$ nombres~?
\end{question}

\begin{question}[$\star\star\star$ \iqroles{QR, QT, Trader} \iqfirm{hedge fund, teneur de marché} \iqlevel{confirmé}]\label{iq:iv1:expectation-and-indicators:22}
Dix cartes numérotées de 1 à 10 sont distribuées au hasard sur dix cases numérotées de 1 à 10, une carte par case. Chaque carte
qui tombe sur son propre numéro est retirée, avec sa case. Les cartes restantes sont battues et distribuées à nouveau
sur les cases restantes, et l'opération se répète jusqu'à ce qu'il ne reste plus aucune carte. Quelle est l'espérance du nombre de distributions~?
\end{question}

\begin{question}[$\star\star\star\star$ \iqroles{QR, QT} \iqfirm{teneur de marché, hedge fund} \iqlevel{confirmé}]\label{iq:iv1:expectation-and-indicators:23}
Vous tirez des nombres aléatoires indépendants, chacun uniformément réparti entre 0 et 1, et vous continuez tant que chaque nombre est
plus grand que celui qui précède. Vous vous arrêtez au premier nombre plus petit que celui qui précède. Quelle est l'espérance
de ce dernier nombre~?
\end{question}

\begin{question}[$\star\star\star\star$ \iqroles{QR, QT} \iqfirm{teneur de marché, hedge fund, crypto} \iqlevel{confirmé}]\label{iq:iv1:expectation-and-indicators:24}
Vous lancez un dé équilibré jusqu'à ce que chaque face soit apparue au moins une fois. Lorsque vous vous arrêtez, certaines faces sont
apparues exactement une fois~; la dernière face nouvelle l'a toujours été. Quelle est l'espérance du nombre de faces apparues exactement une fois~?
\end{question}

\begin{question}[$\star\star\star\star$ \iqroles{QR, QT} \iqfirm{hedge fund, teneur de marché} \iqlevel{confirmé}]\label{iq:iv1:expectation-and-indicators:25}
Vous lancez un dé équilibré jusqu'à ce que chaque face soit apparue au moins deux fois. Mettez en place un calcul exact qu'un recruteur
pourrait suivre au tableau, et donnez l'espérance du nombre de lancers à deux décimales.
\end{question}
